\documentclass[10pt, aspectratio=169]{beamer}
\usepackage[theorembox=color]{adenc-beamer}
% Options: theorembox=none|plain|color, hideproofs, hidemarkings

\newcommand{\cmd}[1]{\texttt{\detokenize{#1}}}

\title{\texttt{adenc-beamer.sty} Feature Showcase}
\author{Aden Chen}
\date{\today}


\begin{document}


\begin{frame}
  \titlepage
\end{frame}


\begin{frame}{Contents}
  \tableofcontents
\end{frame}


% ================================================
\section{Package Options}
% ================================================

\begin{transitionframe}
  \begin{center}{\Huge Package Options}\end{center}
\end{transitionframe}

\begin{frame}{Package Options}
  Pass options via \cmd{\usepackage[opts]{adenc-beamer}}.
  \begin{itemize}
    \item \cmd{theorembox=none|plain|color} --- theorem box style
      (this slide uses \cmd{theorembox=color})
    \item \cmd{hideproofs} --- hide all \texttt{proof} environments
    \item \cmd{hidemarkings} --- hide \cmd{\markabove}/\cmd{\markbelow} annotations
  \end{itemize}

  \bigskip
  Beamer-specific defaults: \cmd{theorembox} defaults to \cmd{plain}
  (boxed, no color) rather than \cmd{none}.
\end{frame}


% ================================================
\section{Theorem Environments}
% ================================================

\begin{transitionframe}
  \begin{center}{\Huge Theorem Environments}\end{center}
\end{transitionframe}

\begin{frame}{Theorem Environments}
  \twocolumns{0.5}{0.5}{%
    \begin{definition}
      A \vocab{group} is a set $G$ with a binary operation satisfying
      closure, associativity, identity, and inverses.
    \end{definition}
    \begin{theorem}[Lagrange]
      If $H \leq G$ is a subgroup of a finite group $G$, then
      $\lvert H \rvert$ divides $\lvert G \rvert$.
    \end{theorem}
    \begin{lemma}
      Left cosets of $H$ partition $G$.
    \end{lemma}
  }{%
    \begin{corollary}
      The order of any element divides $\lvert G \rvert$.
    \end{corollary}
    \begin{example}
      $\ZZ/6\ZZ$ has subgroups of orders $1, 2, 3, 6$.
    \end{example}
    \begin{remark}
      The converse of Lagrange's theorem fails in general.
    \end{remark}
  }
\end{frame}

\begin{frame}{More Environments}
  \twocolumns{0.5}{0.5}{%
    \begin{proposition}
      Every cyclic group is abelian.
    \end{proposition}
    \begin{claim}
      $S_3$ is the smallest non-abelian group.
    \end{claim}
    \begin{conjecture}
      Every finite simple group is either cyclic of prime order
      or an alternating group $A_n$ ($n \geq 5$), or a group of Lie type, \ldots
    \end{conjecture}
  }{%
    \begin{problem}
      Classify all groups of order $p^2$.
    \end{problem}
    \begin{note}
      These environments all have starred (\texttt{*}) unnumbered variants,
      e.g.\ \cmd{theorem*}, \cmd{definition*}.
    \end{note}
    \begin{proof}
      Omit. \qed
    \end{proof}
  }
\end{frame}


% ================================================
\section{Math Commands}
% ================================================

\begin{transitionframe}
  \begin{center}{\Huge Math Commands}\end{center}
\end{transitionframe}

\begin{frame}{Number Sets \& Blackboard Bold}
  \twocolumns{0.5}{0.5}{%
    \textbf{Number sets (long and short forms):}
    \[
      \NN,\ \ZZ,\ \QQ,\ \RR,\ \CC,\ \FF,\ \PP
    \]
    equivalent to \cmd{\bN, \bZ, \bQ, \bR, \bC, \bF, \bP}.

    \bigskip
    \textbf{Mathcal shortcuts} \cmd{\cA}--\cmd{\cZ}:
    \[
      \cA,\ \cB,\ \cF,\ \cG,\ \cH,\ \cL,\ \cM,\ \cX
    \]
  }{%
    \textbf{Mathscr shortcuts} \cmd{\sA}--\cmd{\sZ}:
    \[
      \sA,\ \sB,\ \sF,\ \sL,\ \sM,\ \sX
    \]

    \bigskip
    \textbf{Imaginary unit:} $\ii$ \quad (\cmd{\ii})

    \bigskip
    \textbf{Indicator function:} $\ind{A}$ \quad (\cmd{\ind{A}})
  }
\end{frame}

\begin{frame}{Delimiters \& Operators}
  \twocolumns{0.5}{0.5}{%
    \textbf{Auto-sizing delimiters:}
    \[
      \abs*{\frac{x}{y}},\quad \norm*{\frac{v}{w}},\quad
      \ceil*{\frac{n}{2}},\quad \floor*{\frac{n}{2}}
    \]

    \bigskip
    \textbf{Vectors \& matrices:}
    \[
      \vec{v},\quad
      \bmat{1 & 0 \\ 0 & 1},\quad
      \pmat{a & b \\ c & d}
    \]
  }{%
    \textbf{Calculus:}
    \[
      \DD f,\quad \int_0^1 f(x)\,\dd x,\quad
      \int e^{\ii t}\,\dd t
    \]

    \bigskip
    \textbf{Probability:}
    \[
      \E[X],\quad \Var(X),\quad \Cov(X,Y),\quad X \indep Y
    \]
  }
\end{frame}


% ================================================
\section{Annotations \& Markup}
% ================================================

\begin{transitionframe}
  \begin{center}{\Huge Annotations \& Markup}\end{center}
\end{transitionframe}

\begin{frame}{Draft Annotations}
  \textbf{\cmd{\draftnote[Label][color]{text}}} --- inline draft comment.

  \bigskip
  \begin{itemize}
    \item \cmd{\draftnote{fix this}} produces \draftnote{fix this}
    \item \cmd{\draftnote[TODO][blue]{rewrite}} produces \draftnote[TODO][blue]{rewrite}
  \end{itemize}

  \bigskip\bigskip
  \textbf{\cmd{\vocab{word}}} --- bold vocabulary term, as in:
  a \vocab{normal subgroup} $N \trianglelefteq G$ satisfies $gNg^{-1} = N$.
\end{frame}

\begin{frame}{Two-Column Layout}
  \cmd{\twocolumns{lw}{rw}{left content}{right content}} splits a frame:
  \bigskip

  \twocolumns{0.45}{0.45}{%
    \textbf{Left column} (\texttt{lw=0.45})\\
    Content here spans 45\% of \cmd{\textwidth}.
    \begin{itemize}
      \item Item one
      \item Item two
    \end{itemize}
  }{%
    \textbf{Right column} (\texttt{rw=0.45})\\
    Content here spans another 45\%.
    \begin{enumerate}[label=(\roman*)]
      \item First
      \item Second
    \end{enumerate}
  }
\end{frame}


\end{document}
